% Part: first-order-logic
% Chapter: introduction
% Section: sentences

\documentclass[../../../include/open-logic-section]{subfiles}

\begin{document}

\olfileid{fol}{int}{snt}

\olsection{\usetoken{P}{sentence}}

ठीक आहे, आता आपल्याकडे !!{structure}s आणि !!{formula}s यांच्यासाठी
पूर्तिच्या (``यात सत्य'') व्याख्येचा आराखडा आहे. पण तिला हा अतिरिक्त
भाग---!!a{variable} चे मूल्यांकन---आवश्यक आहे आणि आपल्याला हवी होती ती
!!{sentence}s ची व्याख्या. मूल्यांकने कशी काढून टाकायची आणि
!!{sentence}s म्हणजे काय?

आकारिक तर्कशास्त्राशी तुमची पहिली ओळख झाली तेव्हा !!{variable}s आणि
संख्यापक यांच्यातील संबंधाची केलेली चर्चा तुम्हाला बहुधा आठवत असेल.
संख्यापकानंतर नेहमी !!a{variable} येते आणि मग त्या संख्यापकाचा
!!{sentence} मधील ज्या भागाला उपयोग होतो (त्याची ``व्याप्ती''), त्या
भागात ते !!{variable} त्या संख्यापकाने ``बद्ध'' केले आहे असे आपण
समजतो. !!{formula}s मध्ये प्रत्येक !!{variable} ला जुळणारा संख्यापक
असणे आवश्यक नव्हते आणि जुळणारे संख्यापक नसलेली !!{variable}s ही
``मुक्त'' किंवा ``अबद्ध'' असतात. मुक्त !!{variable}s नसलेली सर्व
!!{formula}s म्हणजे !!{sentence}s असे आपण मानू.

!!a{formula}~$!A$ मधील !!a{variable} ची एखादी आस्थिती केव्हा बद्ध असते,
तिला कोणता संख्यापक बद्ध करतो आणि ती केव्हा मुक्त असते, याची अंतःप्रज्ञ
कल्पना समजणे पुन्हा अवघड नाही. हे तपासण्याची, कदाचित कंस मोजण्याची,
एखादी पद्धत तुम्ही शिकला असाल. पण आपण काटेकोर व्याख्येचा आग्रह धरला
पाहिजे---आणि !!{formula}s ची व्याख्या आपण विगमनाने केली असल्यामुळे,
!!a{formula}~$!A$ मधील !!a{variable}~$x$ च्या मुक्त आणि बद्ध आस्थितींची
व्याख्याही विगमनाने देता येते. उदाहरणार्थ, आपल्या सोप्या भाषेसाठी ती
अशी दिसू शकेल:
\begin{enumerate}
  \item $!A$ हे आण्विक असेल, तर त्यातील $x$ च्या सर्व आस्थिती मुक्त
  असतात (म्हणजे, $\Atom{\Obj P}{x}$ मधील $x$ ची आस्थिती मुक्त असते).
  \item $!A$ हे $\lnot !B$ या रूपाचे असेल, तर $\lnot !B$ मधील~$x$ ची
  आस्थिती तेव्हा आणि केवळ तेव्हाच मुक्त असते, जेव्हा~$!B$ मधील~$x$ ची
  संबंधित आस्थिती मुक्त असते (म्हणजे,~$!B$ मधील चरांच्या मुक्त
  आस्थिती म्हणजे नेमक्या $\lnot !B$ मधील संबंधित आस्थिती होत).
  \item $!A$ हे $(!B \land !C)$ या रूपाचे असेल, तर $(!B \land !C)$
  मधील~$x$ ची आस्थिती तेव्हा आणि केवळ तेव्हाच मुक्त असते, जेव्हा~$!B$
  किंवा~$!C$ मधील~$x$ ची संबंधित आस्थिती मुक्त असते.
  \item $!A$ हे $\lexists[x][!B]$ या रूपाचे असेल, तर $!A$ मधील $x$ ची
  कोणतीही आस्थिती मुक्त नसते; ते $\lexists[y][!B]$ या रूपाचे असेल आणि
  $y$ हे~$x$ पेक्षा वेगळे !!{variable} असेल, तर
  $\lexists[y][!B]$ मधील~$x$ ची आस्थिती तेव्हा आणि केवळ तेव्हाच मुक्त
  असते, जेव्हा~$!B$ मधील~$x$ ची संबंधित आस्थिती मुक्त असते.
\end{enumerate}

चरांच्या मुक्त आणि बद्ध आस्थितींची काटेकोर व्याख्या मिळाल्यानंतर आपण
सरळ असे म्हणू शकतो: मुक्त !!{variable}s ची आस्थिती नसलेले कोणतेही
!!{formula} म्हणजे !!a{sentence} होय.

\end{document}
